\begin{table}[htbp]
\centering
\caption{The round-price price-impact premium in large trades}
\label{tab:dimp}
\small
\begin{threeparttable}
\begin{tabular}{@{}lrrrrrrr@{}}
\toprule
Horizon & Side & Pooled & SE & $t$ & Stock-day mean & Stock-days & Ohta \\
\midrule
1s & buy & nan & (nan) & nan & 0.254 & 1,244 & -- \\
1s & sell & nan & (nan) & nan & 0.999 & 1,479 & -- \\
60s & buy & nan & (nan) & nan & -0.091 & 1,243 & 1.26 \\
60s & sell & nan & (nan) & nan & 1.040 & 1,476 & 1.12 \\
300s & buy & nan & (nan) & nan & -0.412 & 1,239 & -- \\
300s & sell & nan & (nan) & nan & 0.805 & 1,468 & -- \\
\bottomrule
\end{tabular}
\begin{tablenotes}[flushleft]\footnotesize
\item $\Delta\mathit{Imp}$ is the volume-weighted midquote change after trades at a last price digit of zero minus the same quantity for other trades, within large trades on each side, in basis points and signed so that a positive value means the price moved the way the initiator pushed it. The pooled column removes stock and day effects and weights each stock-day by the smaller of its two cell counts; standard errors are clustered by stock and by day. Cells with fewer than five trades are dropped. The final column is Ohta's (2026) 2010--2022 mean at the one-minute horizon. $^{*}p<0.1$, $^{**}p<0.05$, $^{***}p<0.01$.
\end{tablenotes}
\end{threeparttable}
\end{table}
